\documentclass[10pt,a4paper]{amsart}
\usepackage[margin=19mm]{geometry}
\usepackage{amsmath,amssymb,mathtools}
\usepackage[scr=rsfs,cal=euler]{mathalfa}
\usepackage{xfrac}
\usepackage[unicode,hidelinks,bookmarks=false]{hyperref}
\newcommand{\ie}{i.e.\ }
\newcommand{\cf}{cf.\ }
\newcommand{\resp}{respectively\ }
\newcommand{\Subsection}{\subsection}
\providecommand{\nobreakdash}{\nobreak\hbox{-}}
\setlength{\emergencystretch}{2em}
\multlinegap0pt
\usepackage[shortlabels]{enumitem}
\usepackage{amssymb}
\usepackage{xcolor}
\newtheorem{lemma}{Lemma}
\newcommand{\fb}{\textnormal{FB}}
\title{The Length of Functional Batch and PIR Codes}
\author{Altan B. K\i l\i\c{c}, Alberto Ravagnani, Flavio Salizzoni}
\date{Source: arXiv:2508.02586v2 (CC BY 4.0)}
\begin{document}
\maketitle

\noindent\emph{Source and licence.} The Length of Functional Batch and PIR Codes. Version arXiv:2508.02586v2 (CC BY 4.0), distributed by arXiv under the Creative Commons Attribution 4.0 International licence, http://creativecommons.org/licenses/by/4.0/, as recorded in the arXiv licence metadata for this exact version and shown on the abstract page https://arxiv.org/abs/2508.02586v2. Full licence notice: \texttt{licenses/arxiv-2508.02586-CC-BY-4.0.txt}.\par\smallskip

\noindent\textit{Editorial scope.} Counted unit: the article's lemma lem:upper\_bound\_t=2 (source0.tex lines 666--669). The article's complete original proof of it is delivered with it (source0.tex lines 670--745, one \texttt{proof} environment of 76 lines). No mathematics is written, repaired or re-derived: the statement and the proof are the article's own lines, in the article's own notation, and no retained line is substituted or re-flowed.  No further source-local context is needed: every cross-reference of every retained line resolves inside the two delivered spans.  Nothing was omitted from any retained span, so the package carries no omission record.  No retained line contains a citation, so the wrapper carries no bibliography and no bibliography entry.  No retained line contains a figure environment, an image-inclusion command or drawing code, so no image is delivered and the package declares no asset dependency.  Editorial formatting only, in addition: an AMS \texttt{amsart} preamble that restates the article's own declarations and macros used by the retained lines, namely the environments \texttt{lemma} on separate counters, together with the standard packages those lines need.  The retained environments print in this document as Lemma 1 in order of appearance, whereas the article prints the same items with the numbering of its own full document.  The complete native source of the article as distributed in the arXiv e-print for this exact version is bundled unchanged as \texttt{sources/182677/source0.tex}.
\begin{lemma}
\label{lem:upper_bound_t=2}
$\fb(k,2,2) \le \left\lceil \frac{3k}{2} \right\rceil.$
\end{lemma}

\begin{proof}
Let $m$ be a positive integer. The lemma is equivalent to show that $\fb(2m,2,2) \le 3m$ and $\fb(2m+1,2,2) \le 3m+2$. Consider the matrices
$$M_{\textnormal{even}} = \left( I_{2m} \mid r_1 \ldots r_m\right) \mbox{ and }M_{\textnormal{odd}} = \left( I_{2m+1} \mid r'_1 \ldots r'_m e_1\right),$$

where~$r_i = e_{2i-1}+e_{2i}$ and $r'_i = e_{2i}+e_{2i+1}$ for $i \in \{1,\ldots,m\}$. Observe that the sum of the columns in both matrices equal to the zero vector, proving that any vector can be served twice. Thus, we assume that $L =\{a,b\}$ is the list of vectors to be served where $a \neq b$. 

\textbf{Case 1: }Let $k=2m$. We will show that $M_{\textnormal{even}}$ can be used to serve any two vectors proving $\fb(2m,2,2) \le 3m$. Let $p_i = \{2i-1,2i\}$ for $i \in \{1,\ldots,m\}$. Define $$J_i =(|\sigma(a) \cap p_i|,|\sigma(b) \cap p_i|,|\sigma(a) \cap \sigma(b) \cap p_i|)$$ for $i \in \{1,\ldots,m\}$.
Fix $i \in \{1,\ldots,m\}$. We explain the recovery scheme in Table \ref{table:rec1}. This simply follows from the construction of the matrix $M_{\textnormal{even}}$ and the fact that the columns in the non-systematic part have pairwise disjoint supports. Thus, for any of the $m$ intervals of size one needs at most 3 columns. Thus $\fb(2m,2,2) \le 3m$.

\textbf{Case 2: }Let $k=2m+1$. We will show that $M_{\textnormal{odd}}$ can be used to serve any two vectors, proving that $\fb(2m+1,2,2) \le 3m+2$. Let $p'_i = \{2i,2i+1\}$ for $i \in \{1,\ldots,m\}$. Define $$J_i =(|\sigma(a) \cap p_i|,|\sigma(b) \cap p_i|,|\sigma(a) \cap \sigma(b) \cap p_i|)$$ for $i \in \{1,\ldots,m\}$. Fix $i \in \{1,\ldots,m\}$. We explain the recovery scheme in Table \ref{table:rec2}. This simply follows from the construction of the matrix $M_{\textnormal{odd}}$ and the fact that the columns in the non-systematic part have pairwise disjoint supports. Thus, for any of the $m$ intervals of size one needs at most 3 columns. The only thing left is to check the first coordinates. There, it is possible that $a_1=b_1=1$, forcing one to use $e_1$ both in the systematic and the non-sytematic part of $M_{\textnormal{odd}}$. Therefore, $\fb(2m+1,2,2) \le 3m+2$. \qedhere


\begin{table}
\parbox{.49\linewidth}{
\centering
\resizebox{0.5\textwidth}{!}{
\begin{tabular}{|c| c|}
\hline
$J_i$& number of columns used to recover $L$\\
\hline
 $(0,0,0)$  & none. \\
\hline
$(0,1,0)$  & 1 systematic col.    \\
\hline
$(0,2,0)$  & 1 non-systematic col.  \\
\hline
$(1,0,0)$  & 1 systematic col.  \\
\hline
$(1,1,0)$  & 2 systematic col.   \\
\hline
$(1,1,1)$  & 2 systematic col. and 1 non-systematic col.   \\
\hline
$(1,2,1)$  & 1 systematic col. and 1 non-systematic col.   \\
\hline
$(2,0,0)$  & 1 non-systematic col.   \\
\hline
$(2,1,1)$  & 1 systematic col. and 1 non-systematic col.   \\
\hline
$(2,2,2)$  & 2 systematic col. and 1 non-systematic col.  \\
\hline
\end{tabular}}
\caption{$k=2m$ is even.} \label{table:rec1}
}
\hfill
\parbox{.49\linewidth}{
\centering
\resizebox{0.5\textwidth}{!}{
\begin{tabular}{|c| c|}
\hline
$J_i$& number of columns used to recover $L$\\
\hline
 $(0,0,0)$  & none. \\
\hline
$(0,1,0)$  & 1 systematic col.    \\
\hline
$(0,2,0)$  & 1 non-systematic col.  \\
\hline
$(1,0,0)$  & 1 systematic col.  \\
\hline
$(1,1,0)$  & 2 systematic col.   \\
\hline
$(1,1,1)$  & 2 systematic col. and 1 non-systematic col.   \\
\hline
$(1,2,1)$  & 1 systematic col. and 1 non-systematic col.   \\
\hline
$(2,0,0)$  & 1 non-systematic col.   \\
\hline
$(2,1,1)$  & 1 systematic col. and 1 non-systematic col.   \\
\hline
$(2,2,2)$  & 2 systematic col. and 1 non-systematic col.  \\
\hline
\end{tabular}}
\caption{$k=2m+1$ is odd.} \label{table:rec2}
}
\end{table} 
\end{proof}

\end{document}
